\documentclass[conference]{IEEEtran}
\IEEEoverridecommandlockouts
\usepackage{cite}
\usepackage{amsmath,amssymb,amsfonts}
\usepackage{graphicx}
\usepackage{textcomp}
\usepackage{xcolor}
\usepackage[utf8]{inputenc}
\usepackage[spanish,es-nodecimaldot]{babel}
\addto\captionsspanish{\renewcommand{\tablename}{Tabla}}
\addto\captionsspanish{\renewcommand{\figurename}{Fig.}}
\usepackage{hyperref}
\usepackage{booktabs}
\usepackage{url}
\usepackage{listings}
\usepackage{tikz}
\usetikzlibrary{shapes.geometric, arrows.meta, positioning, fit, backgrounds}

\lstset{
  basicstyle=\ttfamily\scriptsize,
  breaklines=true,
  columns=fullflexible,
  frame=single,
  backgroundcolor=\color{black!3},
  numbers=left,
  numberstyle=\tiny\color{gray},
  showstringspaces=false,
  captionpos=b,
  xleftmargin=1.2em,
  framexleftmargin=1em
}

\begin{document}

\title{Prueba de esfuerzo de un sistema Asterisk:\\
configuraci\'on y aplicaci\'on de \textit{Load Testing} sobre AGI--TAXI\\
{\footnotesize \textsuperscript{*}Entrega 2 --- Seminario de Voz IP, Universidad de Antioquia, 2026}
}

\author{
\IEEEauthorblockN{Ana Mar\'ia Garz\'on}
\IEEEauthorblockA{\textit{Facultad de Ingenier\'ia} \\
\textit{Universidad de Antioquia}\\
Medell\'in, Colombia}
\and
\IEEEauthorblockN{Juan Esteban Cardozo Rivera}
\IEEEauthorblockA{\textit{Facultad de Ingenier\'ia} \\
\textit{Universidad de Antioquia}\\
Medell\'in, Colombia \\
juan.cardozor@udea.edu.co}
\and
\IEEEauthorblockN{V\'ictor Manuel Restrepo}
\IEEEauthorblockA{\textit{Facultad de Ingenier\'ia} \\
\textit{Universidad de Antioquia}\\
Medell\'in, Colombia}
}

\maketitle

\begin{abstract}
La capacidad operativa de un PBX basado en Asterisk se ve limitada por
el procesamiento de se\~nalizaci\'on SIP, el manejo de canales RTP y la
l\'ogica de negocio invocada por las extensiones. Este documento ---
segunda entrega del proyecto de prueba de esfuerzo del Seminario de Voz
IP --- describe la configuraci\'on aplicada sobre el sistema
AGI--TAXI (Asterisk 18+ sobre Debian) para habilitar una prueba de
carga (\textit{Load Testing}) con SIPp, y reporta las evidencias de su
ejecuci\'on. Se construyeron dos escenarios SIP/RTP id\'enticos en
se\~nalizaci\'on pero diferenciados en procesamiento --- uno con
\texttt{Echo()} y otro con la l\'ogica AGI real de consulta a SQLite ---
para aislar el costo del plano de negocio frente al costo puro de la
pila SIP. La concurrencia se incrementa por escalones de 10 a 300
llamadas, midiendo CPU, memoria, canales activos, jitter y p\'erdida de
paquetes. Los resultados permiten ubicar el punto de inflexi\'on de la
VM y dejar trazabilidad reproducible mediante archivos de
configuraci\'on, escenarios XML y scripts de monitoreo publicados.
\end{abstract}

\begin{IEEEkeywords}
Asterisk, SIPp, VoIP, load testing, stress testing, PJSIP, AGI
\end{IEEEkeywords}

\section{Introducci\'on}

En la Entrega 1 se describieron las principales t\'ecnicas de
\textit{stress testing} aplicables a un PBX Asterisk --- \textit{Load
Testing}, \textit{Stress Testing}, \textit{Spike Testing}, \textit{Soak
Testing} y \textit{Codec Stress Testing} ---, las m\'etricas relevantes
(CPU, CPS, llamadas concurrentes, jitter, p\'erdida de paquetes,
latencia y MOS) y las herramientas para llevarlas a cabo, con SIPp
\cite{ref_sipp} como n\'ucleo del proceso. La presente entrega traslada
ese marco te\'orico a un sistema real, el proyecto AGI--TAXI
\cite{ref_agitaxi}, y aplica una prueba de carga sobre \'el, dejando
evidencias reproducibles de la configuraci\'on aplicada y de los
indicadores capturados.

El objetivo espec\'ifico de esta entrega es responder, sobre la VM del
proyecto, dos preguntas operativas: \textit{(i) cu\'antas llamadas
concurrentes puede sostener el PBX manteniendo un consumo de CPU
inferior al 80\,\% y un porcentaje de llamadas exitosas superior al
99\,\%}, y \textit{(ii) cu\'al es el sobrecosto introducido por la
l\'ogica AGI de negocio (consulta SQLite) frente a una extensi\'on
trivial de eco}. La elecci\'on de \textit{Load Testing} sobre las otras
t\'ecnicas se justifica porque permite curvas monotonas, comparables
entre escenarios, sin riesgo de colgar la VM, y porque produce
evidencias gr\'aficas claras a partir de un \'unico CSV de monitoreo.

\section{Sistema bajo prueba}

\subsection{Arquitectura AGI--TAXI}

El sistema bajo prueba es el desarrollado en el Seminario de Voz IP:
una VM Debian con Asterisk 18+ y un script AGI en Python que consulta
una base SQLite local con datos de demostraci\'on de taxis por c\'odigo
postal. Los componentes relevantes para esta prueba son:

\begin{itemize}
    \item \textbf{Asterisk 18+} con m\'odulo PJSIP, escuchando en
    \texttt{192.168.56.10:5060/UDP}.
    \item \textbf{Endpoints PJSIP} \texttt{1001} y \texttt{1002}
    definidos en \texttt{pjsip.conf} para validaci\'on con softphones.
    \item \textbf{Extensiones del dialplan}: \texttt{100} (servicio
    AGI), \texttt{600} (echo) y \texttt{1001}/\texttt{1002}
    (intercomunicaci\'on interna), todas en el contexto
    \texttt{from-internal}.
    \item \textbf{Script AGI} \texttt{taxi\_agi.py} en
    \texttt{/var/lib/asterisk/agi-bin/}, que abre una conexi\'on
    SQLite a una base sembrada con cinco c\'odigos postales de
    Medell\'in.
\end{itemize}

\subsection{Configuraci\'on a\~nadida para la prueba}

Para evitar interferir con el sistema productivo del proyecto, las
adiciones se aislaron en dos archivos \texttt{\#include} dentro de
\texttt{/etc/asterisk}:

\begin{itemize}
    \item \texttt{pjsip-stress.conf}: declara el endpoint
    \texttt{6000} con autenticaci\'on digest, AOR con
    \texttt{max\_contacts=200} y \texttt{context=from-stress}. El
    endpoint reutiliza los templates \texttt{endpoint-template},
    \texttt{auth-template} y \texttt{aor-template} ya presentes en el
    repositorio del proyecto.
    \item \texttt{extensions-stress.conf}: define el contexto
    \texttt{from-stress} con dos extensiones de carga ---
    \texttt{7000} (Answer + Echo + Hangup) y \texttt{7100} (Answer +
    AGI + Hangup) --- que solo son alcanzables desde el endpoint de
    pruebas.
\end{itemize}

La Fig.~\ref{fig:arquitectura} muestra el montaje final.

\begin{figure}[!t]
\centering
\resizebox{\columnwidth}{!}{%
\begin{tikzpicture}[
  font=\scriptsize,
  every node/.style={align=center},
  vm/.style={rectangle, draw, rounded corners, minimum width=2.4cm, minimum height=1.4cm, thick, fill=blue!8},
  svc/.style={rectangle, draw, rounded corners, minimum width=1.8cm, minimum height=0.7cm, fill=orange!12},
  arr/.style={-{Stealth[length=3mm,width=2.5mm]}, very thick},
  biarr/.style={{Stealth[length=3mm,width=2.5mm]}-{Stealth[length=3mm,width=2.5mm]}, very thick}
]
  \node[vm] (cli) at (0,0)  {VM Cliente\\(SIPp)};
  \node[vm] (srv) at (6,0)  {VM Asterisk\\(AGI--TAXI)};
  \node[svc] (sip) at (6,-1.0) {PJSIP 5060/UDP};
  \node[svc] (rtp) at (6,-1.8) {RTP 10000--20000/UDP};
  \node[svc] (agi) at (6,-2.6) {AGI taxi\_agi.py};
  \node[svc] (db)  at (6,-3.4) {SQLite local};

  \draw[biarr] (cli) -- node[above,font=\tiny]{INVITE/200/BYE} (srv);
  \draw[biarr] (cli) -- ++(0,-1.6) -| node[pos=0.25, below, font=\tiny]{RTP G.711} (rtp);
  \draw[arr]   (srv) -- (sip);
  \draw[arr]   (sip) -- (rtp);
  \draw[arr]   (rtp) -- (agi);
  \draw[arr]   (agi) -- (db);
\end{tikzpicture}
}
\caption{Montaje de la prueba de carga: una VM cliente genera carga SIP/RTP hacia la VM Asterisk del proyecto AGI--TAXI.}
\label{fig:arquitectura}
\end{figure}

\section{Dise\~no experimental}

\subsection{Tipo de prueba}

Se aplica \textit{Load Testing}: se incrementa progresivamente la
concurrencia hasta el l\'imite operativo esperado, manteniendo cada
escal\'on durante un tiempo suficiente para que las m\'etricas se
estabilicen. A diferencia del \textit{Stress Testing}, no se busca el
punto de quiebre catastr\'ofico sino caracterizar la zona de
operaci\'on aceptable y el inicio de la degradaci\'on. Cada escal\'on
se ejecuta dos veces, una con el escenario \texttt{Echo} (extensi\'on
\texttt{7000}) y otra con el escenario \texttt{AGI} (extensi\'on
\texttt{7100}), para aislar el costo del plano de negocio.

\subsection{Escalones de concurrencia y tasa}

Se definen siete escalones (Tabla~\ref{tab:escalones}) con duraci\'on
fija de 120\,s por escal\'on y enfriamiento de 10\,s entre escalones.
La tasa de llegada se mantiene proporcional a la concurrencia objetivo
para acortar el transitorio de rampa.

\begin{table}[!h]
\caption{Escalones del Load Test.}
\label{tab:escalones}
\renewcommand{\arraystretch}{1.2}
\setlength{\tabcolsep}{6pt}
\begin{center}
\footnotesize
\begin{tabular}{@{}rrrr@{}}
\toprule
\textbf{Paso} & \textbf{Concurrencia (\texttt{-l})} & \textbf{CPS (\texttt{-r})} & \textbf{Llamadas (\texttt{-m})} \\
\midrule
1 & 10  & 5  & 600   \\
2 & 25  & 10 & 1\,200 \\
3 & 50  & 15 & 1\,800 \\
4 & 100 & 20 & 2\,400 \\
5 & 150 & 25 & 3\,000 \\
6 & 200 & 30 & 3\,600 \\
7 & 300 & 40 & 4\,800 \\
\bottomrule
\end{tabular}
\end{center}
\end{table}

\subsection{M\'etricas y criterios de aceptaci\'on}

Para cada escal\'on se registran las m\'etricas presentadas en la
Entrega 1: CPU (\%), memoria en uso (MB), llamadas concurrentes
activas en Asterisk, tasa de llamadas exitosas, retransmisiones,
jitter y p\'erdida de paquetes. Se considera un escal\'on como
\textit{aceptable} cuando se cumplen simult\'aneamente:

\begin{itemize}
    \item CPU sostenida (\textit{usr+sys+iow}) inferior al 80\,\%.
    \item Llamadas exitosas $\geq$ 99\,\% de las intentadas.
    \item Jitter promedio $<$ 30\,ms y p\'erdida $<$ 1\,\%.
\end{itemize}

El \'ultimo escal\'on que cumple los tres criterios se reporta como la
\textit{capacidad sostenida} de la VM bajo el escenario probado.

\section{Configuraci\'on aplicada en Asterisk}

\subsection{Endpoint dedicado a la prueba}

El Listing~\ref{lst:pjsip} muestra el contenido del archivo
\texttt{pjsip-stress.conf} incluido al final de
\texttt{/etc/asterisk/pjsip.conf}.

\lstinputlisting[
  language={},
  caption={\texttt{pjsip-stress.conf}: endpoint SIPp con auth digest.},
  label=lst:pjsip
]{../asterisk/pjsip-stress.conf}

\subsection{Contexto de dialplan de carga}

El Listing~\ref{lst:dialplan} muestra el dialplan a\~nadido. La
extensi\'on \texttt{7000} hace s\'olo \texttt{Answer()} + \texttt{Echo()}
y la \texttt{7100} invoca el mismo AGI de producci\'on
\texttt{taxi\_agi.py}, lo que permite comparar el costo unitario por
canal con y sin l\'ogica de negocio.

\lstinputlisting[
  language={},
  caption={\texttt{extensions-stress.conf}: contexto \texttt{from-stress}.},
  label=lst:dialplan
]{../asterisk/extensions-stress.conf}

\subsection{Recarga sin reinicio}

Tras copiar los archivos se ejecuta el script
\texttt{apply-config.sh} que a\~nade los \texttt{\#include} si no
existen y recarga Asterisk con \texttt{dialplan reload} y
\texttt{pjsip reload}, sin caer servicios productivos.

\section{Herramienta SIPp y escenarios}

Se utiliza SIPp 3.6.1 \cite{ref_sipp} sobre una VM cliente
independiente (puede ser la misma host-only del laboratorio). El
escenario base se compone de dos pares INVITE/200OK/ACK --- el primero
para responder al reto digest (\texttt{401 Unauthorized}) y el segundo
con la cabecera \texttt{[authentication]} ya resuelta --- seguido de
una pausa que mantiene la llamada activa y un BYE final. El
Listing~\ref{lst:sipp} reproduce el escenario \textit{sin RTP},
suficiente para estresar el plano de se\~nalizaci\'on. Para estresar
tambi\'en el plano RTP se proporciona la variante
\texttt{uac-echo-rtp.xml} que reproduce un PCAP G.711 contra la
extensi\'on \texttt{Echo()}.

\lstinputlisting[
  language=XML,
  caption={\texttt{uac-echo-no-rtp.xml}: escenario UAC con digest.},
  label=lst:sipp
]{../sipp/uac-echo-no-rtp.xml}

\section{Procedimiento de ejecuci\'on}

El procedimiento completo, reproducible, se sintetiza en cinco pasos:

\begin{enumerate}
    \item \textbf{Preparaci\'on de la VM Asterisk}: aplicar
    configuraci\'on con \texttt{sudo ./apply-config.sh} y verificar el
    endpoint con \texttt{asterisk -rx 'pjsip show endpoint 6000'}.
    \item \textbf{Preparaci\'on de la VM cliente}: instalar SIPp con
    \texttt{sudo ./install-tools.sh}.
    \item \textbf{Arranque del monitoreo} en la VM Asterisk:
    \texttt{./monitor.sh results/monitor.csv} en una sesi\'on aparte.
    \item \textbf{Ejecuci\'on del Load Test} en la VM cliente:
    \texttt{./run-load.sh}. El script orquesta los siete escalones de
    la Tabla~\ref{tab:escalones}, deja CSV con estad\'isticas por
    escal\'on y logs de errores y pantalla de SIPp.
    \item \textbf{An\'alisis}: \texttt{python3 plot.py
    results/monitor.csv results/} genera las gr\'aficas y la tabla se
    consolida en la Secci\'on~\ref{sec:resultados}.
\end{enumerate}

\section{Resultados y evidencias}
\label{sec:resultados}

\subsection{Resumen por escal\'on}

La Tabla~\ref{tab:resultados} concentra los indicadores capturados por
escal\'on para el escenario \texttt{Echo}. La \'ultima fila marca el
escal\'on aceptado seg\'un los criterios de la Secci\'on anterior.

\begin{table}[!h]
\caption{Resultados del Load Test (escenario Echo). Los valores deben actualizarse tras la ejecuci\'on real con los datos capturados por \texttt{monitor.sh} y \texttt{run-load.sh}.}
\label{tab:resultados}
\renewcommand{\arraystretch}{1.2}
\setlength{\tabcolsep}{4pt}
\begin{center}
\footnotesize
\begin{tabular}{@{}rrrrrrr@{}}
\toprule
\textbf{Conc.} & \textbf{CPS} & \textbf{CPU\,(\%)} & \textbf{Mem\,(MB)} & \textbf{OK\,(\%)} & \textbf{Jitter\,(ms)} & \textbf{Loss\,(\%)} \\
\midrule
10   & 5   & --- & --- & --- & --- & --- \\
25   & 10  & --- & --- & --- & --- & --- \\
50   & 15  & --- & --- & --- & --- & --- \\
100  & 20  & --- & --- & --- & --- & --- \\
150  & 25  & --- & --- & --- & --- & --- \\
200  & 30  & --- & --- & --- & --- & --- \\
300  & 40  & --- & --- & --- & --- & --- \\
\bottomrule
\end{tabular}
\end{center}
\end{table}

\subsection{Curvas capturadas}

La Fig.~\ref{fig:cpu} muestra la curva CPU vs canales activos generada
por \texttt{plot.py} a partir del CSV de \texttt{monitor.sh}. La
Fig.~\ref{fig:net} muestra el trafico SIP/RTP en kbps. Ambas im\'agenes
deben reemplazarse por las generadas en la ejecuci\'on real.

\begin{figure}[!h]
\centerline{\includegraphics[width=0.95\columnwidth]{placeholder_cpu}}
\caption{CPU (\%) y canales activos de Asterisk durante el Load Test.}
\label{fig:cpu}
\end{figure}

\begin{figure}[!h]
\centerline{\includegraphics[width=0.95\columnwidth]{placeholder_net}}
\caption{Tr\'afico SIP/RTP (RX/TX en kbps) durante el Load Test.}
\label{fig:net}
\end{figure}

\subsection{Evidencias adicionales}

Como evidencia complementaria se anexan en el repositorio del proyecto:
(i) el archivo \texttt{summary.csv} producido por \texttt{run-load.sh}
con totales por escal\'on; (ii) los \texttt{screen.log} y
\texttt{error.log} de SIPp por escal\'on; (iii) la salida de
\texttt{asterisk -rx 'core show channels count'} y
\texttt{asterisk -rx 'pjsip show endpoints'} antes, durante y al
finalizar la prueba; y (iv) capturas de pantalla de \texttt{htop} y
\texttt{sngrep} durante el escal\'on m\'aximo aceptado.

\subsection{Comparativa Echo vs AGI}

Repetir los escalones cambiando \texttt{SERVICE=7100} en
\texttt{run-load.sh} permite observar el sobrecosto del AGI Python
sobre SQLite. Se espera una reducci\'on de la capacidad sostenida y un
aumento del consumo de CPU en estado de usuario --- atribuible al
intervalo de espera por entrada DTMF y a la consulta SQLite --- pero
no en el plano RTP, ya que la se\~nalizaci\'on y los flujos de medios
son id\'enticos.

\section{Discusi\'on}

La prueba ejecutada permite caracterizar el sistema sobre la
infraestructura realmente disponible al proyecto, sin requerir
hardware adicional. La separaci\'on de la configuraci\'on de carga en
archivos \texttt{\#include} preserva la integridad del repositorio
AGI--TAXI y facilita revertir el escenario de pruebas (basta con
remover las dos l\'ineas \texttt{\#include} y recargar Asterisk).
Asimismo, el uso de un endpoint dedicado con autenticaci\'on digest
documenta una buena pr\'actica frente al hosting an\'onimo que muchas
gu\'ias de SIPp sugieren para simplificar.

Como limitaciones, la concurrencia m\'axima alcanzable depende del
sizing de la VM (vCPUs, RAM, NIC virtual). El procedimiento descrito
queda parametrizado para que cualquier ejecuci\'on en un host m\'as
potente pueda extender los escalones m\'as all\'a de 300, simplemente
ajustando el arreglo \texttt{LEVELS} de \texttt{run-load.sh}.

\section{Conclusiones}

La Entrega 2 cierra el ciclo planteado en la Entrega 1: las t\'ecnicas
y herramientas descritas se materializan en archivos de configuraci\'on,
escenarios XML y scripts ejecutables sobre la VM real del proyecto
AGI--TAXI. La prueba de carga aplicada sobre las extensiones
\texttt{7000} (Echo) y \texttt{7100} (AGI) entrega como salidas la
capacidad sostenida del PBX, el sobrecosto del plano de negocio y un
CSV reproducible que sirve como l\'inea base para futuras pruebas de
\textit{Stress}, \textit{Spike} o \textit{Soak} sobre el mismo
sistema.

\begin{thebibliography}{00}
\bibitem{ref_sipp} SIPp Developers, ``SIPp Documentation (v3.6.1),''
2023. [En l\'inea]. Disponible: \url{https://sipp.readthedocs.io/}.
[Accedido: 2026-05-31].
\bibitem{ref_asterisk} Sangoma Technologies, ``Asterisk 18
Documentation,'' 2024. [En l\'inea]. Disponible:
\url{https://docs.asterisk.org/}. [Accedido: 2026-05-31].
\bibitem{ref_pjsip} Sangoma Technologies, ``PJSIP Configuration
Reference,'' 2024. [En l\'inea]. Disponible:
\url{https://docs.asterisk.org/Configuration/Channel-Drivers/SIP/Configuring-res_pjsip/}.
[Accedido: 2026-05-31].
\bibitem{ref_agitaxi} J. C. Restrepo, ``AGI--TAXI: Sistema de consulta
de disponibilidad de taxi por voz,'' Universidad de Antioquia,
Repositorio p\'ublico, 2026. [En l\'inea]. Disponible:
\url{https://github.com/JuanC101195/AGI-TAXI}.
\bibitem{ref_davidson} J. Davidson, J. Peters, M. Bhatia,
S. Kalidindi y S. Mukherjee, \textit{Voice over IP Fundamentals},
2.\textsuperscript{a} ed. Indianapolis: Cisco Press, 2006.
\end{thebibliography}

\end{document}
